% =============================================================================
% exercises/ejercicios_cap05.tex
% Ejercicios propuestos — Capítulo 5: La Atmósfera
% =============================================================================

\section*{Ejercicios Propuestos}
\addcontentsline{toc}{section}{Ejercicios propuestos}
\label{sec:ejercicios_cap05}

\begin{bloque_ejercicios}[Ejercicios --- Cap\'itulo 4: La Atm\'osfera]

\begin{ejercicios}

% ----- Ejercicio 1: Estructura y capas de la atmósfera ----------------------
\item Con base en la descripci\'on de las capas atmosf\'ericas
del cap\'itulo, complete la siguiente tabla:

{%
\centering
\begin{tabular}{lp{2.5cm}p{2.5cm}p{3.0cm}}
\toprule
\textbf{Capa} & \textbf{Altitud (km)} & \textbf{Temp. (K)} &
\textbf{Caracter\'istica principal} \\
\midrule
Troposfera    & & & \\
Estratosfera  & & & \\
Mesosfera     & & & \\{0.5em}
Term\'osfera  & & & \\
Exosfera      & & & \\
\bottomrule
\end{tabular}
\par
}%

\begin{enumerate}[label=\alph*)]
    \item ?`En cu\'al de las capas se producen los fen\'omenos
          meteorol\'ogicos cotidianos (lluvia, viento, tormentas)?
          ?`Por qu\'e?
    \item ?`En qu\'e capa se encuentra la mayor concentraci\'on
          de ozono estratosf\'erico? ?`A qu\'e altitud aproximada?
    \item La term\'osfera tiene temperaturas de hasta
          \SI{2000}{\kelvin}, pero no se percibe ese calor.
          Explique por qu\'e en t\'erminos de la densidad molecular.
\end{enumerate}

\textit{[Respuesta parcial: Troposfera: 0--17~km, 200--300~K,
mezclado r\'apido; Estratosfera: 17--50~km, contiene la capa de
ozono; Mesosfera: 50--100~km, 130--250~K;
Term\'osfera: 100--500~km, 1000--2000~K;
Exosfera: $>$500~km, escape al espacio]}

% ----- Ejercicio 2: Composición de la atmósfera ----------------------------
\item Con base en el \textbf{Cuadro~\ref{Atmcomp}}:

\begin{enumerate}[label=\alph*)]
    \item Los cuatro gases m\'as abundantes en la atm\'osfera son
          \ce{N2}, \ce{O2}, Ar y \ce{CO2}. Exprese sus fracciones
          molares en porcentaje (\%) y en partes por mill\'on en
          volumen (ppm$_v$).

{%
\centering
\begin{tabular}{lccl}
\toprule
\textbf{Gas} & \textbf{Fracci\'on molar} & \textbf{\%} & \textbf{ppm$_v$} \\
\midrule
\ce{N2}  & 0.78     & & \\
\ce{O2}  & 0.21     & & \\
Ar       & 0.0093   & & \\
\ce{CO2} & $390\times10^{-6}$ & & \\
\bottomrule
\end{tabular}
\par
}%

    \item La concentraci\'on de \ce{CO2} en el texto se indica
          como 410~ppm (valor actual), mientras que la tabla
          muestra 390~ppm (dato de 1972). Calcule el incremento
          absoluto (ppm) y relativo (\%) entre ambos valores.

    \item El metano (\ce{CH4}) se encuentra a $1.7\times10^{-6}$
          mol/mol. Expr\'eselo en ppb$_v$.
          ?`Cu\'antas veces m\'as abundante es el \ce{CO2}
          que el \ce{CH4}?
\end{enumerate}

\textit{[Respuesta:
(a) \ce{N2}: 78\%, 780~000~ppm; \ce{O2}: 21\%, 210~000~ppm;
    Ar: 0.93\%, 9~300~ppm; \ce{CO2}: 0.039\%, 390~ppm;
(b) $\Delta$ = 20~ppm; incremento = $20/390 \times 100 = 5.1\%$;
(c) 1.7~ppb$_v$; \ce{CO2}/\ce{CH4} $\approx$ 229 veces]}

% ----- Ejercicio 3: El ozono estratosférico --------------------------------
\item El ozono estratosf\'erico se forma y destruye mediante
el ciclo de Chapman:
\[
    \ce{O2 ->[$h\nu$] 2O^.} \quad
    \ce{O^. + O2 ->  O3} \quad
    \ce{O3 ->[$h\nu$] O2 + O^.}
\]

\begin{enumerate}[label=\alph*)]
    \item A partir de la \textbf{Figura~\ref{atmoO3}}, indique
          a qu\'e altitud se alcanza la concentraci\'on m\'axima
          de ozono y estime su valor aproximado en
          mol\'eculas/cm$^3$.
    \item Los CFC son inertes en la troposfera pero se descomponen
          en la estratosfera. Explique por qu\'e la radiaci\'on UV
          que los descompone no llega a la troposfera.
    \item Un \'atomo de Cl liberado por un CFC puede destruir
          hasta $10^5$ mol\'eculas de \ce{O3} de forma catal\'itica.
          Escriba el ciclo catal\'itico b\'asico:
          \ce{Cl + O3 ->} \underline{\hspace{2cm}} y
          \ce{ClO + O^. ->} \underline{\hspace{2cm}}.
    \item ?`Por qu\'e el agujero de ozono se forma especialmente
          en el Polo Sur durante la primavera austral y no de forma
          uniforme en todo el planeta?
\end{enumerate}

\textit{[Respuesta: (c) \ce{Cl + O3 -> ClO + O2};
\ce{ClO + O^. -> Cl + O2};
neto: \ce{O3 + O^. -> 2O2}]}

% ----- Ejercicio 4: Efecto invernadero y temperatura del planeta -----------
\item El efecto invernadero natural mantiene la temperatura
promedio del planeta en \SI{15}{\celsius} en lugar de
\SI{-18}{\celsius}.

\begin{enumerate}[label=\alph*)]
    \item Calcule la diferencia de temperatura ($\Delta T$) que
          aporta el efecto invernadero natural y expr\'esela
          tambi\'en en Kelvin.
    \item Los principales gases de efecto invernadero son
          \ce{CO2}, \ce{CH4}, \ce{N2O} y \ce{H2O}. ?`Cu\'al de
          ellos es el m\'as abundante en la atm\'osfera seg\'un
          el cap\'itulo? ?`Y el que tiene mayor PCG-100 a\~nos
          de los tres primeros, seg\'un el \textbf{Cuadro~\ref{WGP}}?
    \item Explique por qu\'e el vapor de agua (\ce{H2O}) no se
          incluye en el \textbf{Cuadro~\ref{WGP}} de potenciales
          de calentamiento global, a pesar de ser el GEI m\'as
          abundante.
    \item Las actividades humanas han incrementado la concentraci\'on
          de GEI. Mencione tres fuentes antropog\'enicas descritas
          en el cap\'itulo y el gas principal que emite cada una.
\end{enumerate}

\textit{[Respuesta:
(a) $\Delta T = 15 - (-18) = \SI{33}{\celsius} = \SI{33}{\kelvin}$;
(b) el m\'as abundante es \ce{H2O} (0--4\%); el mayor PCG-100 es
    \ce{N2O} con 298;
(d) quema de combustibles f\'osiles (\ce{CO2}),
    ganader\'ia (\ce{CH4}), fertilizantes (\ce{N2O})]}

% ----- Ejercicio 5: Potencial de calentamiento global (PCG) ----------------
\item Con los datos del \textbf{Cuadro~\ref{WGP}}:

\begin{enumerate}[label=\alph*)]
    \item El PCG expresa cu\'antas veces m\'as calentamiento
          produce 1~kg de un gas respecto a 1~kg de \ce{CO2}
          en un horizonte temporal dado. Si se emiten
          \SI{1000}{\kilo\gram} de \ce{CH4} a la atm\'osfera,
          ?`a cu\'antos kg equivalentes de \ce{CO2} corresponde
          en el horizonte de 100~a\~nos? ?`Y en 20~a\~nos?

    \item Calcule las emisiones equivalentes de \ce{CO2}
          (en kg-\ce{CO2}eq) para:
          \begin{enumerate}[label=\roman*)]
              \item \SI{500}{\kilo\gram} de \ce{N2O}
                    (horizonte 100 a\~nos)
              \item \SI{10}{\kilo\gram} de CFC-12
                    (horizonte 100 a\~nos)
              \item \SI{1}{\kilo\gram} de \ce{SF6}
                    (horizonte 100 a\~nos)
          \end{enumerate}

    \item Ordene los siguientes gases de mayor a menor PCG-100
          y comente la relaci\'on entre la vida del compuesto y
          su potencial de calentamiento:
          \ce{CH4}, \ce{N2O}, HFC-23, \ce{SF6}, PFC-14.

    \item El bromuro de metilo (\ce{CH3Br}) tiene una vida de
          solo 0.7~a\~nos. Sin embargo, su PCG-20 (17) es mayor
          que su PCG-100 (5). Explique esta relaci\'on.
\end{enumerate}

\textit{[Respuesta:
(a) 100~a\~nos: $1000 \times 25 = 25\,000$~kg-\ce{CO2}eq;
    20~a\~nos: $1000 \times 72 = 72\,000$~kg-\ce{CO2}eq;
(b-i) $500 \times 298 = 149\,000$~kg-\ce{CO2}eq;
    (b-ii) $10 \times 10\,900 = 109\,000$~kg-\ce{CO2}eq;
    (b-iii) $1 \times 22\,800 = 22\,800$~kg-\ce{CO2}eq]}

% ----- Ejercicio 6: Señales del cambio climático ---------------------------
\item Las se\~nales del cambio clim\'atico descritas en el
cap\'itulo son observables en distintos sistemas del planeta.

\begin{enumerate}[label=\alph*)]
    \item Para cada se\~nal, identifique el mecanismo f\'isico
          o qu\'imico que la explica:

{%
\centering
\begin{tabular}{lp{6cm}}
\toprule
\textbf{Se\~nal} & \textbf{Mecanismo} \\
\midrule
Incremento de \ce{CO2} atmosf\'erico & \\[0.4em]
Calentamiento de oc\'eanos            & \\[0.4em]
Reducci\'on de glaciares              & \\[0.4em]
Blanqueamiento de corales             & \\[0.4em]
Aumento del nivel del mar             & \\
\bottomrule
\end{tabular}
\par
}%

    \item Explique la diferencia entre \textit{clima} y
          \textit{tiempo meteorol\'ogico} con un ejemplo concreto
          de cada uno, usando la definici\'on del cap\'itulo.

    \item Las emisiones continuas de GEI podr\'ian aumentar la
          temperatura global entre \SIrange{1}{5}{\celsius} para
          el a\~no 2100. ?`Cu\'ales son dos consecuencias espec\'ificas
          sobre los recursos costeros mencionadas en el cap\'itulo?
\end{enumerate}

% ----- Ejercicio 7: Ciclos de Milankovitch ---------------------------------
\item Los ciclos de la Tierra descritos en la Secci\'on~\ref{ctierra}
modulan la radiaci\'on solar recibida en la superficie.

\begin{enumerate}[label=\alph*)]
    \item Identifique los tres ciclos de Milankovitch mencionados
          en el cap\'itulo, indicando su nombre y periodo en a\~nos:

{%
\centering
\begin{tabular}{lll}
\toprule
\textbf{Par\'ametro} & \textbf{Nombre del ciclo} & \textbf{Per\'iodo (a\~nos)} \\
\midrule
Bamboleo del eje   & & \\
Inclinaci\'on del eje & & \\
Variaci\'on orbital  & & (dos valores) \\
\bottomrule
\end{tabular}
\par
}%

    \item Durante los per\'iodos glaciales los niveles de \ce{CO2}
          disminuyen significativamente. ?`Qu\'e relaci\'on existe
          entre los ciclos orbitales y los niveles de \ce{CO2}
          seg\'un el cap\'itulo?

    \item A diferencia de los ciclos naturales, el calentamiento
          actual ocurre en d\'ecadas, no en miles de a\~nos.
          ?`Qu\'e implica esta diferencia en t\'erminos de la
          capacidad de adaptaci\'on de los ecosistemas?
\end{enumerate}

% ----- Ejercicio 8: Forzamiento radiativo ----------------------------------
\item El forzamiento radiativo (FR) cuantifica el desbalance
energ\'etico en la tropopausa, en unidades de \si{\watt\per\square\metre}.

\begin{enumerate}[label=\alph*)]
    \item ?`Qu\'e significa un FR positivo? ?`Y un FR negativo?
          Proporcione un ejemplo de cada caso con compuestos
          del \textbf{Cuadro~\ref{WGP}} o de la figura de
          forzantes radiativos.

    \item El \ce{CO2} tiene un PCG-100 = 1 (es el gas de
          referencia). Sin embargo, tiene el mayor forzamiento
          radiativo total acumulado. ?`C\'omo es posible si
          su PCG es el m\'as bajo de todos? Explique en t\'erminos
          de \textbf{concentraci\'on y masa emitida}.

    \item El \ce{SF6} tiene una vida de 3~200~a\~nos y un
          PCG-100 de 22~800. Si hoy se emite 1~kg de \ce{SF6},
          ?`durante cu\'antas generaciones humanas (asumiendo
          25~a\~nos/generaci\'on) permanecer\'a activo en la
          atm\'osfera?

    \item Los aerosoles de sulfato tienen un FR negativo
          (efecto enfriante). Explique por qu\'e su presencia
          en la estratosfera puede enfriar el clima temporalmente
          (como despu\'es de erupciones volc\'anicas).
\end{enumerate}

\textit{[Respuesta parcial:
(c) $3\,200/25 = 128$~generaciones humanas]}

% ----- Ejercicio 9: Protocolo de Montreal y CFCs ---------------------------
\item El Protocolo de Montreal (1987) regula las sustancias que
agotan la capa de ozono listadas en el \textbf{Cuadro~\ref{WGP}}.

\begin{enumerate}[label=\alph*)]
    \item Identifique las tres sustancias controladas por el
          Protocolo de Montreal con mayor PCG-100 en el cuadro.
          Ord\'enelas de mayor a menor potencial.

    \item El CFC-12 (\ce{CCl2F2}) tiene una vida de 100 a\~nos
          y PCG-100 de 10~900. El HCFC-123 (\ce{CHCl2CF3})
          tiene una vida de 1.3 a\~nos y PCG-100 de 77. ?`Por qu\'e
          los HCFC fueron considerados sustitutos intermedios
          de los CFC? Mencione ventaja e inconveniente.

    \item Un pa\'is emite \SI{500}{\tonne} de CFC-11
          (\ce{CCl3F}) al a\~no. Calcule las emisiones equivalentes
          en \si{\tonne} de \ce{CO2} usando el PCG-100.

    \item Los HFC fueron introducidos como sustitutos de los
          HCFC por no agotar el ozono. Sin embargo, aparecen
          en el cuadro con PCG muy elevados. ?`Qu\'e dilema
          ambiental representa esto para la pol\'itica
          clim\'atica global?
\end{enumerate}

\textit{[Respuesta:
(a) Halon-11301 (8~480) $>$ CFC-11 (6~730) \ldots
    (revisar tabla para orden completo);
(c) $500 \times 6\,730 = 3\,365\,000$~t-\ce{CO2}eq]}

% ----- Ejercicio 10 (avanzado): Cálculo integrador -------------------------
\item \textbf{[Avanzado]} Integraci\'on de conceptos atmosf\'ericos.

\begin{enumerate}[label=\alph*)]
    \item La densidad del aire a nivel del mar es de
          $\approx \SI{1.225}{\kilo\gram\per\cubic\metre}$
          y la presi\'on atmosf\'erica es de \SI{101325}{\pascal}.
          Usando la ley del gas ideal, estime el n\'umero de
          mol\'eculas por cm$^3$ a \SI{25}{\celsius}.
          \textit{Use: $n/V = P/(RT)$,
          $R = \SI{8.314}{\joule\per\mole\per\kelvin}$,
          y $N_A = 6.022\times10^{23}$~mol\'ec/mol.}

    \item Con la fracci\'on molar del \ce{N2} (0.78), calcule
          la presi\'on parcial del nitr\'ogeno en la atm\'osfera
          a nivel del mar en Pa y en atm\'osferas.

    \item La concentraci\'on de \ce{CH4} es $1.7\times10^{-6}$
          mol/mol. A partir del n\'umero total de mol\'eculas/cm$^3$
          obtenido en (a), calcule el n\'umero de mol\'eculas de
          \ce{CH4} por cm$^3$.

    \item Si la concentraci\'on de \ce{CO2} aumenta de 390 a
          410~ppm en el periodo 1972--2025, y asumiendo un
          incremento lineal, estime la tasa anual de aumento
          en ppm/a\~no. ?`Es consistente este valor con los
          registros actuales del observatorio Mauna Loa
          ($\approx$~2--3~ppm/a\~no)?
\end{enumerate}

\textit{[Respuesta:
(a) $n/V = 101325/(8.314\times298.15)
= 40.9$~mol/m$^3$;
mol\'ec/cm$^3 = 40.9\times10^{-6}\times6.022\times10^{23}
= 2.46\times10^{19}$~mol\'ec/cm$^3$;
(b) $P_{\ce{N2}} = 0.78\times101325
= \SI{79034}{\pascal} = \SI{0.780}{\atm}$;
(c) $1.7\times10^{-6}\times2.46\times10^{19}
= 4.2\times10^{13}$~mol\'ec-\ce{CH4}/cm$^3$;
(d) $20~\text{ppm}/53~\text{a\~nos}
\approx 0.38$~ppm/a\~no --- menor que el valor actual;
el incremento es acelerado (no lineal)]}

\end{ejercicios}

\end{bloque_ejercicios}
% =============================================================================
% FIN DE ejercicios_cap05.tex
% =============================================================================
